\subsection{Mục tiêu và Phạm vi dự án}

\subsubsection{Mục tiêu tổng quan}

Mục tiêu tổng quan của dự án
là nghiên cứu, thiết kế và xây dựng
một hệ thống trợ lý pháp luật
thông minh, ứng dụng công nghệ trí tuệ nhân tạo (AI)
kết hợp cùng nền tảng kiến trúc vi dịch vụ (Microservices).
Hệ thống được phát triển nhằm giải quyết
những khó khăn trong việc tra cứu thủ công,
qua đó hỗ trợ người dùng phổ thông và
các chuyên gia pháp lý tiếp cận, truy vấn và
giải đáp các vấn đề liên quan đến
hệ thống văn bản pháp luật
Việt Nam một cách nhanh chóng, chính xác.
Đặc biệt, dự án hướng tới việc xây dựng một trợ lý AI
không chỉ cung cấp câu trả lời cuối cùng
mà còn đưa ra các bước lập luận logic,
minh bạch, kèm theo trích dẫn
nguồn văn bản pháp lý chính thống,
giúp giảm thiểu tối đa hiện tượng \enquote{ảo giác AI} (hallucination) thường gặp trong các mô hình ngôn ngữ lớn.

\subsubsection{Mục tiêu cụ thể}

Để đạt được mục tiêu tổng quan, dự án hướng tới hoàn thành các mục tiêu cụ thể
về mặt tính năng và trải nghiệm người dùng như sau:

\begin{itemize}
\item \textbf{Đảm bảo an toàn thông tin và định danh tài khoản}:
Xây dựng quản lý tài khoản,
tích hợp các phương thức xác thực đa yếu tố
nhằm bảo vệ tối đa dữ liệu
và quyền riêng tư của người dùng.

\item \textbf{Nâng cao tính minh bạch và độ tin cậy trong tư vấn pháp lý}:
Xử lý các truy vấn pháp lý phức tạp
bằng cách không chỉ đưa ra câu trả lời cuối cùng,
mà còn phải diễn giải chi tiết quá trình suy luận
và cung cấp chính xác các nguồn tài liệu tham khảo.

\item \textbf{Quy trình thanh toán tự động}:
Tích hợp các cổng thanh toán
để tự động hóa việc quản lý các gói dịch vụ thuê bao,
xử lý giao dịch
mang lại trải nghiệm nâng cấp dịch vụ an toàn, tự động.

\item \textbf{Xử lý tài liệu người dùng}:
Phát triển dịch vụ tiếp nhận
và xử lý đa dạng các định dạng văn bản được tải lên,
tự động trích xuất thông tin, vector hóa tài liệu
một cách nhanh chóng, chính xác và bảo mật.

\item \textbf{Tối ưu hóa khả năng tiếp cận và tương tác}:
Xóa bỏ rào cản nhập liệu
bằng cách cung cấp trải nghiệm đàm thoại rảnh tay,
cho phép người dùng tương tác trực tiếp với hệ thống qua giọng nói tự nhiên.

\end{itemize}

\subsubsection{Phạm vi dự án}

Dự án tập trung nghiên cứu, thiết kế
và phát triển các thành phần cốt lõi
trong một hệ thống vi dịch vụ,
giới hạn trong các phạm vi sau:

\begin{itemize}

\item \textbf{Phân hệ quản trị người dùng và bảo mật}:
Cài đặt hệ thống xác thực thông qua Supabase,
hỗ trợ đăng nhập Email/mật khẩu,
Google OAuth và xác thực hai yếu tố (2FA - TOTP).

\item \textbf{Phân hệ Thu thập dữ liệu pháp luật}:
Thiết lập đường ống dữ liệu (Data Pipeline)
tự động thu thập và đồng bộ hóa văn bản quy phạm pháp luật
từ hai nguồn là \url{https://vbpl.vn}
và CSDL Pháp Điển (\url{https://phapdien.moj.gov.vn}),
có áp dụng thuật toán mã băm để bỏ qua tài liệu không có thay đổi nội dung.

\item \textbf{Phân hệ trợ lý AI và quản lý trò chuyện}:
Phát triển dịch vụ trò chuyện (NestJS) và dịch vụ chatbot (Python) giao tiếp qua gRPC và RabbitMQ.
Ứng dụng kỹ thuật RAG (Retrieval-Augmented Generation) và LLM để trả lời theo cơ chế streaming, xuất luồng suy luận (reasoning details).
Giới hạn các tính năng quản lý ở mức: tạo thư mục, ghi chú, đánh dấu (bookmark) và chia sẻ liên kết trò chuyện.

\item \textbf{Phân hệ xử lý tài liệu cá nhân}:
Quản lý xử lý các tài liệu do người dùng tải lên thông qua luồng tác vụ bất đồng bộ (Celery worker).
Tải tệp từ Cloudflare R2,
trích xuất cấu trúc và nội dung sang Markdown bằng thư viện \texttt{Docling}, phân đoạn văn bản bằng \texttt{RecursiveCharacterTextSplitter}, tạo vector biểu diễn ngữ nghĩa bằng mô hình \texttt{keepitreal/vietnamese-sbert} và lưu trữ vào CSDL vector Qdrant để phục vụ truy xuất. Áp dụng cơ chế kiểm tra trùng lặp (Deduplication) để tối ưu hiệu năng hệ thống.

\item \textbf{Phân hệ giao tiếp bằng giọng nói (Voice AI)}:
Tích hợp nền tảng LiveKit để thu phát âm thanh thời gian thực. Ứng dụng các giải pháp Text-to-Speech bao gồm edge-tts, piper-tts và elevenlabs để tổng hợp giọng nói phản hồi.

\item \textbf{Phân hệ quản lý thanh toán nội địa}:
Xây dựng một dịch vụ thanh toán độc lập triển khai trên Heroku,
sử dụng BullMQ
tự động
xử lý hủy các giao dịch thanh toán quá hạn nếu hết thời gian chờ
và
giao tiếp bất đồng bộ qua Kafka.
Hỗ trợ quản lý trạng thái
thuê bao VIP (Subscriptions)
qua các cổng thanh toán nội địa Việt Nam bao gồm: SePay, VNPAY, ZaloPay và MoMo.
\end{itemize}
